\chapter{Results}
\label{chap:05}
\setlength{\parskip}{1em}

\section{Machine Learning Model Evaluation}

\subsection{Lead Scoring Model Performance}

Model comparison was performed on the held-out 20\% test partition
(n~=~1,486~records). \Cref{tab:lead_results} reports the final test-set
metrics for all candidate models.

\begin{table}[H]
  \centering
  \caption{Lead scoring model comparison on held-out test set}
  \label{tab:lead_results}
%  \setstretch{1.1}
  \begin{tabularx}{\textwidth}{@{}L{3.8cm} C{1.5cm} C{1.5cm} C{1.5cm}
    C{1.5cm} C{1.6cm} C{1.6cm}@{}}
    \toprule
    \thead{Model} & \thead{Acc.} & \thead{Prec.} & \thead{Rec.} &
    \thead{F1} & \thead{AUC-ROC} & \thead{PR-AUC} \\
    \midrule
    Logistic Regression (baseline) &
      0.774 & 0.701 & 0.698 & 0.699 & 0.801 & 0.712 \\
    Random Forest &
      0.801 & 0.742 & 0.756 & 0.749 & 0.851 & 0.779 \\
    \textbf{Gradient Boosting} (selected) &
      \textbf{0.823} & \textbf{0.798} & \textbf{0.863} &
      \textbf{0.829} & \textbf{0.891} & \textbf{0.841} \\
    XGBoost &
      0.819 & 0.791 & 0.855 & 0.822 & 0.887 & 0.836 \\
    \bottomrule
  \end{tabularx}
\end{table}

The Gradient Boosting model was selected as the production lead scoring
model, achieving F1 = 0.829 and AUC-ROC = 0.891, both exceeding the
project objectives of F1 $\geq$ 0.78 and AUC-ROC $\geq$ 0.80. The five
highest-importance features by mean absolute SHAP value were:
\texttt{days\_since\_last\_contact}, \texttt{activity\_count\_30d},
\texttt{email\_open\_rate}, \texttt{deal\_stage\_max}, and
\texttt{call\_count\_90d}. The dominance of behavioural engagement
features is consistent with the findings of \citet{tsai2010}.

\subsection{Churn Prediction Model Performance}

\begin{table}[H]
  \centering
  \caption{Churn prediction model comparison on held-out test set}
  \label{tab:churn_results}
%  \setstretch{1.1}
  \begin{tabularx}{\textwidth}{@{}L{3.8cm} C{1.5cm} C{1.5cm} C{1.5cm}
    C{1.5cm} C{1.6cm} C{1.6cm}@{}}
    \toprule
    \thead{Model} & \thead{Acc.} & \thead{Prec.} & \thead{Rec.} &
    \thead{F1} & \thead{AUC-ROC} & \thead{PR-AUC} \\
    \midrule
    Logistic Regression (baseline) &
      0.793 & 0.651 & 0.712 & 0.680 & 0.834 & 0.721 \\
    \textbf{Random Forest} (selected) &
      \textbf{0.820} & \textbf{0.689} & \textbf{0.741} &
      \textbf{0.714} & \textbf{0.872} & \textbf{0.768} \\
    Gradient Boosting &
      0.816 & 0.682 & 0.748 & 0.714 & 0.868 & 0.762 \\
    \bottomrule
  \end{tabularx}
\end{table}

The Random Forest classifier was selected as the production churn model
with AUC-ROC = 0.872, exceeding the project objective of AUC-ROC $\geq$
0.80. The most important SHAP features for churn prediction were
\texttt{tenure\_months}, \texttt{monthly\_revenue},
\texttt{days\_since\_last\_contact}, \texttt{support\_tickets\_30d}, and
\texttt{call\_count\_90d}. The primacy of \texttt{tenure\_months} aligns
with the customer lifetime value literature \citep{reichheld2001}.

Customer churn risk was classified into three tiers for dashboard
presentation: Low Risk (churn probability below 0.30), Medium Risk
(0.30--0.60), and High Risk (above 0.60), enabling sales managers to
prioritise retention outreach.

\subsection{SHAP Feature Importance Visualisation}

\Cref{fig:shap_lead} presents a schematic representation of the SHAP
beeswarm summary plot for the lead scoring model. High values of
\texttt{activity\_count\_30d} and low values of
\texttt{days\_since\_last\_contact} are the strongest positive predictors
of lead conversion, consistent with the engagement-centricity finding of
\citet{tsai2010}.

\begin{figure}[H]
  \centering
  \begin{tikzpicture}[
    bar/.style={fill=accentblue!70, draw=primaryblue, minimum height=0.4cm},
    label/.style={font=\footnotesize\ttfamily, anchor=east},
    value/.style={font=\footnotesize, anchor=west}
  ]
    \def\scale{30}  % 1 unit = 30pt, so 0.30 fills 9pt total scale
    % Feature bars: (value, label, y-position)
    \foreach \val/\feat/\yi in {
      0.214/days\_since\_last\_contact/5,
      0.187/activity\_count\_30d/4,
      0.142/email\_open\_rate/3,
      0.118/deal\_stage\_max/2,
      0.093/call\_count\_90d/1
    }{
      \pgfmathsetmacro{\barw}{\val*\scale*10}
      \node[label] at (-0.1, \yi*0.6) {\feat};
      \fill[bar] (0, \yi*0.6-0.15) rectangle (\barw pt, \yi*0.6+0.15);
      \node[value] at (\barw pt + 2pt, \yi*0.6) {\val};
    }
    \draw[->] (0,0.2) -- (0,3.4) node[above, font=\tiny] {};
    \draw[->] (0,0.2) -- (70pt,0.2) node[right, font=\footnotesize]
      {Mean $|\text{SHAP}|$};
    % x-axis ticks
    \foreach \x/\lbl in {0/0.00, 21pt/0.07, 42pt/0.14, 63pt/0.21}{
      \draw (\x, 0.2) -- (\x, 0.12);
      \node[font=\tiny, anchor=north] at (\x, 0.12) {\lbl};
    }
  \end{tikzpicture}
  \caption{Mean absolute SHAP values for the lead scoring model (top 5 features)}
  \label{fig:shap_lead}
\end{figure}

